The live-site parameter can be sharpened when most sites remain immutable.
The distinction is between a threshold which merely exists and a site whose
descendant response has actually changed.

\begin{corollary}[Cactus cycles with bounded productive sites]
\label{cor:aesp-cd-cactus-productive-sites}
Retain the exact-real, static, single-positive-load cactus setting of
\cref{cor:aesp-cd-cactus-live-sites}.  After a cycle block $B$ first closes,
call a cycle site \emph{productive} when its aggregate descendant response
has changed because of a routed singleton admission.  Let $p_B$ be the total
number of distinct productive sites during the remaining trace.  Require the
following immutable-run contract:
\begin{enumerate}[label=(\roman*)]
\item at closure, choose the parent articulation (or the root site) as the
      sole permanent base anchor and compress every other site's current
      snapshot---even one with a finite live threshold---into dormant runs;
\item between the base anchor and productive sites, every dormant cycle site keeps the same local
      diagonal, load, descendant threshold, and active/pinned mark;
\item the first routed change below a dormant site promotes that site before
      applying the change, after which the site is maintained explicitly;
\item every descendant mark or response change is reported on its unique
      cactus route to the root;
\item equality at a threshold remains inactive and ties use a fixed label
      order.
\end{enumerate}
Then the term $(r_B+1)\log(2+L_B)$ in
\cref{eq:aesp-cd-cactus-live-sites-work} may be replaced by
$(p_B+1)\log^2(2+L_B)$.  Equivalently, the cycle contribution is
\begin{equation}
 \widetilde O\!\left(
   \sum_{v\in S^\star}
   \sum_{B\in\mathcal P(v)}(p_B+1)\log^2(2+L_B)
 \right),
 \label{eq:aesp-cd-cactus-productive-sites-work}
\end{equation}
plus the tree and pre-closure term in
\cref{eq:aesp-cd-cactus-live-sites-work}.  In particular, if
$\max_B p_B=\operatorname{polylog}(2+\vol(S^\star))$, the RPPR work remains
\begin{equation*}
 \widetilde O\!\left(
  \frac{1+\vol(S^\star)}{\sqrt\alpha}
 \right),
\end{equation*}
even when the number of dormant live sites is unbounded.

This statement does not cover finite-inner response drift, online factor
revelation inside a dormant run, or an unreported child mutation.  It does not
resolve cactus traces with unboundedly many productive sites.
\end{corollary}

\begin{proof}
At cycle closure, rebuild the static representation with only the stipulated
base anchor explicit; all other frozen sites belong to dormant runs.  Later
cuts are precisely the distinct promoted productive sites together with that
anchor.  On a dormant run
$I$ with endpoint values $(z_L,z_R)$, exact elimination gives, for every site
$j\in I$,
\begin{equation*}
 x_j=A_jz_L+B_jz_R+C_j,
 \qquad A_j,B_j\geq0.
\end{equation*}
The coefficients and the descendant threshold $\theta_j$ are immutable by
hypothesis.  Thus every dormant safety condition is one fixed half-plane
\begin{equation*}
 A_jz_L+B_jz_R+C_j\leq\theta_j.
\end{equation*}
Preprocess these normalized homogeneous half-planes---not their pulled-back
scalar thresholds---in the dual planar hull for every canonical interval of a
balanced segment tree on the cycle order.  A maximal dormant run is a
disjoint union of $O(\log L_B)$ canonical intervals.

Conditioned on the parent articulation scalar $t$, inverse positivity makes
both endpoint values a nonnegative affine ray in $t$.  Exact line/ray
shooting in each static planar hull returns the first strict boundary crossing
and a labeled row in $O(\log L_B)$ work.  Zero ray slope is treated as an
immediate record when the row is already positive and as infinity otherwise;
equality is inactive.  Therefore all dormant runs are queried in
$O((p_B+1)\log^2(2+L_B))$ work.

The base anchor, productive sites, and run endpoints form a cyclic or path Stieltjes
skeleton of order $O(p_B+1)$.  Eliminating each dormant run contributes one
two-port Schur factor, so two tridiagonal solves recover the affine skeleton
response to $t$ in $O(p_B+1)$ work.  On first mutation, splitting the run and
promoting the touched site logically retires every canonical interval
containing that site: no physical deletion is needed, because future
maximal-run covers decompose around it and never use those intervals.  When
$p_B=0$, the wraparound dormant run may have the
same permanent anchor at both ends and is handled by the same half-plane query.
Hence every site is promoted once, and preprocessing, promotions, and final
recovery sum to $\widetilde O(\vol(S^\star))$.

An admission below $B$ changes its aggregate response only at the unique
cycle articulation on the admission-to-root route.  The promotion-before-
mutation rule therefore preserves every dormant hull exactly.  Summing the
per-visit cost over block routes proves
\cref{eq:aesp-cd-cactus-productive-sites-work}.  The support-radius argument
from \cref{cor:aesp-cd-cactus-live-sites} gives the final specialization.
\end{proof}

The bounded-productive-site hypothesis can be removed at the price of a
square-root block factor.  This interpolation is useful when a block has many
mutable sites but is still much larger than the number of sites touched in a
typical response epoch.

\begin{corollary}[Epoch rebuilding for unrestricted productive sites]
\label{cor:aesp-cd-cactus-productive-epochs}
Retain the exact static cactus setting above, but impose no bound on $p_B$.
For every closed cycle block $B$, let $J_B$ be the number of routed singleton
admissions after closure whose descendant-response mutation visits $B$ (with
repeated visits to the same site counted separately).  For any integer
$1\leq k_B\leq L_B$, there is an exact reporter whose post-closure work on
$B$ is
\begin{equation}
 \widetilde O\!\left(
   L_B+J_B\left(k_B+\frac{L_B}{k_B}\right)
 \right).
 \label{eq:aesp-cd-cactus-productive-epoch-work}
\end{equation}
In particular, $k_B=\lceil\sqrt{L_B}\rceil$ gives
$\widetilde O(L_B+J_B\sqrt{L_B})$.  More sharply, if $p_B$ denotes the
number of distinct productive sites over the whole post-closure trace, the
same algorithm automatically gives
\begin{equation}
 \widetilde O\!\left(
  L_B+J_B\min\{p_B+1,\sqrt{L_B}\}
 \right).
 \label{eq:aesp-cd-cactus-productive-hybrid-work}
\end{equation}
Here $p_B$ is only an a posteriori analysis parameter; the online reporter
uses $L_B$ and does not need to know $p_B$.
Including bridges, pre-closure work, and
final recovery gives the unconditional instance-sensitive cactus bound
\begin{equation*}
 \widetilde O\!\left(
  1+\vol(S^\star)+
  \sum_{B\ \mathrm{cycle}}J_B
       \min\{p_B+1,\sqrt{L_B}\}
 +\sum_{v\in S^\star}\operatorname{dist}(o,v)
 \right).
\end{equation*}
Consequently the target RPPR product follows under the alternative route
geometry promise
\begin{equation*}
 \max_{v\in S^\star}
 \sum_{B\in\mathcal P(v)}
       \min\{p_B+1,\sqrt{L_B}\}
 =\widetilde O(1+R_\star).
\end{equation*}
For example, this promise holds when every used parent--child passage through
$B$ traverses
$\widetilde\Omega(\min\{p_B+1,\sqrt{L_B}\})$ edges of that cycle, because
those passage lengths add along a cactus block route.
This is not a universal support-radius product bound.  Consider the accounting
shape consisting of a chain of $k$ length-$L$ cycle blocks, at least
$\lceil\sqrt L\rceil$ distinct productive side sites on every block, and $M$
terminal descendant singleton admissions.  Then $p_B\geq\sqrt L$ and
$J_B\geq M$ on every ancestor block, so the displayed accounting permits
$\Theta(kM\sqrt L)$.  Taking $k=L=m$ and $M=m^2$ gives
$\Theta(m^{7/2})$, whereas the added side volume is lower order and volume
times radius is only $\Theta(m^3)$.  This is a separation of the bound, not an
algorithmic lower bound or a claim that every such ledger is an RPPR
trajectory.
\end{corollary}

\begin{proof}
At the start of an epoch, rebuild in $\widetilde O(L_B)$ work the canonical
static half-plane hulls from the current factors, choose the permanent anchor,
and reset the touched set to empty.  During the epoch, apply the preceding
productive-site construction with at most $k_B$ explicit touched sites.  A
routed response update and its exact next-event query therefore cost
$\widetilde O(k_B)$, including the skeleton solve and dormant-run queries.

Rebuild immediately after the $k_B$-th distinct site is touched.  Every
completed epoch is charged to at least $k_B$ routed admissions, so there are
at most $\lfloor J_B/k_B\rfloor$ completed rebuilds after the initial one.
Repeated mutations at an already touched site cause no extra rebuild and are
already charged by the $J_Bk_B$ term.  Thus the total is
\begin{equation*}
 \widetilde O\!\left(
  L_B+J_Bk_B+\frac{J_B}{k_B}L_B
 \right),
\end{equation*}
which proves \cref{eq:aesp-cd-cactus-productive-epoch-work}.  Summing the
one-time builds, the tree route ledger, and the block terms gives the stated
cactus bound.  If the whole trace has $p_B<k_B$ distinct productive sites,
no rebuild occurs and every visit sees at most $p_B+1$ skeleton sites and
dormant runs.  Otherwise the preceding $k_B$ bound applies.  Taking
$k_B=\lceil\sqrt{L_B}\rceil$ and using
$p_B+1\leq k_B\leq2\sqrt{L_B}$ in the no-rebuild branch proves
\cref{eq:aesp-cd-cactus-productive-hybrid-work}.

Reversing the order of summation gives
\begin{equation*}
 \sum_B J_B\min\{p_B+1,\sqrt{L_B}\}
 \leq
 \sum_{v\in S^\star}\sum_{B\in\mathcal P(v)}
       \min\{p_B+1,\sqrt{L_B}\},
\end{equation*}
which proves the route-geometry specialization.  The final chain ledger gives
the announced limitation.  Moreover
$k_B+L_B/k_B\geq2\sqrt{L_B}$, so the square-root factor is optimal within
this flat strategy of a full block rebuild after one fixed number of distinct
touches.  This observation does not rule out a multilevel persistent hull.
\end{proof}
